\section{Proofs for KLM preferential representation}
\label{annex:p08-klm-preferential-representation}

\begin{annexitem}{annex:p08.klm.1}{thm:p08-klm-basis}{The needed rule derivations, before constructing a model}
\begin{proof}
Reflexivity gives \(\alpha\sim\alpha\); Right Weakening with
\(\alpha\models\beta\) gives Supraclassicality. Conversely
Supraclassicality includes Reflexivity. Thus the only nonexplicit rule
needed to pass from the native basis to the expanded seven is And.
Given \(\alpha\sim\beta,\alpha\sim\gamma\), Cautious Monotonicity
gives \(\alpha\wedge\beta\sim\gamma\). Supraclassicality gives
\(\alpha\wedge\beta\wedge\gamma\sim\beta\wedge\gamma\).
Cut first at antecedent \(\alpha\wedge\beta\), then at antecedent
\(\alpha\), gives \(\alpha\sim\beta\wedge\gamma\).
The converse basis already contains all native rules, so both bases
define the same relations, without Rational Monotony.

If \(\alpha\sim\beta,\beta\sim\alpha,\alpha\sim\gamma\),
Cautious Monotonicity gives \(\alpha\wedge\beta\sim\gamma\).
Left Logical Equivalence changes this to \(\beta\wedge\alpha\sim\gamma\),
and Cut with \(\beta\sim\alpha\) gives \(\beta\sim\gamma\).
Interchanging the antecedents proves equality of their consequence sets.

For the implication rule S, suppose \(\alpha\wedge\beta\sim\gamma\).
Right Weakening gives \(\alpha\wedge\beta\sim\beta\to\gamma\).
Supraclassicality gives
\(\alpha\wedge\neg\beta\sim\beta\to\gamma\).
Or and Left Logical Equivalence, using the classical equivalence of
\((\alpha\wedge\beta)\vee(\alpha\wedge\neg\beta)\) with \(\alpha\),
give \(\alpha\sim\beta\to\gamma\). Conversely, if And and Or
are available, S changes \(\alpha\wedge\beta\sim\gamma\) to
\(\alpha\sim\beta\to\gamma\); And with \(\alpha\sim\beta\),
followed by Right Weakening, derives Cut. This proves the claimed
alternative basis rather than assuming Cut in that direction.

Finally, from \(\alpha\sim\gamma\) and \(\beta\sim\delta\),
Right Weakening gives \(\alpha\sim\gamma\vee\delta\) and
\(\beta\sim\gamma\vee\delta\); Or gives their stated mixed
conclusion. All classical equivalences/entailments used are true at
every valuation, hence in the same fixed \(U\). No uniform substitution
of a background constraint has been performed.
\end{proof}
\end{annexitem}

\begin{annexitem}{annex:p08.klm.2}{thm:p08-klm-soundness}{Each rule and the precise use of smoothness}
\begin{proof}
Reflexivity holds because every minimal \(\alpha\)-state satisfies
\(\alpha\). Left Logical Equivalence holds because equivalent formulas
over \(U\) have identical state sets. Right Weakening holds because
\(l[S]\subseteq U\) and every label satisfying the old consequent
satisfies the new one. These same observations prove Supraclassicality.
And holds because a minimal state satisfying both consequents satisfies
their conjunction.

For Or, \(\widehat{\alpha\vee\beta}=\widehat\alpha\cup\widehat\beta\).
A state minimal in the union is minimal in every one of these two
subsets to which it belongs. The two premise assertions therefore make
each such state satisfy the common consequent, proving Or.

If \(\alpha\sim_W\beta\), every state minimal in
\(\widehat\alpha\) belongs to \(\widehat{\alpha\wedge\beta}\)
and is minimal there, since it is a subset of \(\widehat\alpha\).
This proves one inclusion between the minimal sets. For the converse,
let \(s\) be minimal in \(\widehat{\alpha\wedge\beta}\).
If it is not minimal in \(\widehat\alpha\), smoothness supplies
\(t\prec s\) minimal in \(\widehat\alpha\). The premise assertion
makes \(t\) satisfy \(\beta\), so it lies in
\(\widehat{\alpha\wedge\beta}\), contradicting minimality of \(s\).
Thus the minimal sets are equal, including the empty-set cases.
The equality proves Cautious Monotonicity and the stated preservation
of all consequences. For Cut a weaker argument suffices: every minimal
\(\alpha\)-state already lies in, and is minimal in, the conjunction
set, so the second Cut premise gives its consequent. This last argument
does not use smoothness. Hence every primitive and derived sheet rule
has been checked individually under its actual hypotheses.
\end{proof}
\end{annexitem}

\begin{annexitem}{annex:p08.klm.3}{thm:p08-klm-normal-witness}{Finite satisfiability and the infinite normal theory}
\begin{proof}
If \(\alpha\sim\beta\), the definition of \(N(\alpha)\) makes
every member satisfy \(\beta\). For the converse suppose
\(\alpha\not\sim\beta\), and put
\(T=C(\alpha)\cup\{\neg\beta\}\).
If a finite subset of \(T\) had no model in \(U\), there would
be finitely many \(\delta_1,\ldots,\delta_n\in C(\alpha)\) such
that \(\bigwedge_i\delta_i\models\beta\).
Repeated And gives \(\alpha\sim\bigwedge_i\delta_i\), and Right
Weakening then gives \(\alpha\sim\beta\), a contradiction.
For \(n=0\) the empty conjunction is \(\top\), and
\(\alpha\sim\top\) follows from Reflexivity and Right Weakening;
the same argument applies.
Thus every finite subset of \(T\) has a model in \(U\).
The explicitly assumed compactness supplies \(m\models T\), which
is \(\alpha\)-normal and falsifies \(\beta\). This proves both
directions even when \(C(\alpha)\) is infinite. In particular it
does not define a normal world by merely naming one, nor select a
representative world for every formula simultaneously.
\end{proof}
\end{annexitem}

\begin{annexitem}{annex:p08.klm.4}{thm:p08-klm-preorder}{All preorder and transport dependencies}
\begin{proof}
Reflexivity of \(\leq\) follows from Reflexivity of \(\sim\) and
\(\alpha\vee\alpha\leftrightarrow\alpha\). The identities
\(\alpha\vee\beta\leq\alpha,\beta\) follow the same way.

Suppose \(\alpha\leq\beta\leq\gamma\), and put
\(D=\alpha\vee\beta\vee\gamma\). Mixed Or applied to
\(\alpha\vee\beta\sim\alpha\) and
\(\beta\vee\gamma\sim\beta\) gives
\(D\sim\alpha\vee\beta\).
Supraclassicality gives \(\alpha\vee\beta\sim D\). The mutual
antecedent-equivalence rule therefore transfers
\(\alpha\vee\beta\sim\alpha\) to \(D\sim\alpha\).
Mixed Or applied to \(\alpha\vee\beta\sim\alpha\) and
\(\gamma\sim\gamma\) gives \(D\sim\alpha\vee\gamma\).
Cautious Monotonicity followed by Left Logical Equivalence yields
\(\alpha\vee\gamma\sim\alpha\), which is
\(\alpha\leq\gamma\). This proves transitivity syntactically.

For the first transport, let \(\alpha\leq\beta\) and
\(\beta\sim\delta\). Put \(E=\alpha\vee\beta\).
Since \(E\wedge\beta\) is equivalent to \(\beta\), S gives
\(E\sim\beta\to\delta\). As \(E\sim\alpha\), Cautious
Monotonicity and Left Logical Equivalence give
\(\alpha\sim\beta\to\delta\). An \(\alpha\)-normal world
satisfying \(\beta\) consequently satisfies every such \(\delta\),
so is \(\beta\)-normal.

For the second transport under \(\alpha\leq\beta\leq\gamma\),
\(D\wedge(\beta\vee\gamma)\) is equivalent to \(\beta\vee\gamma\).
Thus S applied to \(\beta\vee\gamma\sim\beta\) gives
\(D\sim(\beta\vee\gamma)\to\beta\), and Right Weakening gives
\(D\sim\gamma\to\beta\). We already established
\(D\sim\alpha\vee\beta\). Cautious Monotonicity first restricts
the antecedent to \(\alpha\vee\beta\), then, using
\(\alpha\vee\beta\sim\alpha\), to \(\alpha\). Therefore
\(\alpha\sim\gamma\to\beta\). An \(\alpha\)-normal world
satisfying \(\gamma\) satisfies \(\beta\); the first transport
then makes it \(\beta\)-normal. These derivations supply the
source's intermediate rules rather than relying on semantic completeness
that is still to be proved.
\end{proof}
\end{annexitem}

\begin{annexitem}{annex:p08.klm.5}{thm:p08-klm-canonical-order}{Strictness of the actual copied-state order}
\begin{proof}
If \((m,\alpha)\prec(m,\alpha)\), its definition would give
\(m\not\models\alpha\). But \(m\in N(\alpha)\) and
\(\alpha\sim\alpha\), so \(m\models\alpha\), a contradiction.
For transitivity, suppose
\((m_0,\alpha_0)\prec(m_1,\alpha_1)\prec(m_2,\alpha_2)\).
Then \(\alpha_0\leq\alpha_1\leq\alpha_2\) and
\(m_0\not\models\alpha_1\). Preorder transitivity gives
\(\alpha_0\leq\alpha_2\). If \(m_0\models\alpha_2\), the
second normal-world transport would make it satisfy \(\alpha_1\),
a contradiction. Thus \(m_0\not\models\alpha_2\), exactly the
remaining guard required for
\((m_0,\alpha_0)\prec(m_2,\alpha_2)\).
This proves transitivity rather than replacing the valuation guard by
preorder transitivity alone.
\end{proof}
\end{annexitem}

\begin{annexitem}{annex:p08.klm.6}{thm:p08-klm-canonical-minima}{Both minimal-state implications}
\begin{proof}
Suppose \(m\models\alpha\) and \(\beta\leq\alpha\).
If an \(\alpha\)-state \((n,\gamma)\) preceded \((m,\beta)\),
then \(\gamma\leq\beta\leq\alpha\), \(n\in N(\gamma)\),
\(n\models\alpha\), and \(n\not\models\beta\).
Normal-world transport forces \(n\models\beta\), a contradiction.
Thus \((m,\beta)\) is minimal in \(\widehat\alpha\).

Conversely, suppose it is minimal there. Then \(m\models\alpha\).
If \(\beta\not\leq\alpha\), then
\(\alpha\vee\beta\not\sim\beta\). The normal-world witness
gives \(n\in N(\alpha\vee\beta)\) with \(n\not\models\beta\).
It satisfies \(\alpha\vee\beta\), hence \(\alpha\).
Also \(\alpha\vee\beta\leq\beta\). Therefore
\((n,\alpha\vee\beta)\prec(m,\beta)\) is an
\(\alpha\)-predecessor, contradicting minimality. This establishes
\(\beta\leq\alpha\), and the characterization is complete.
\end{proof}
\end{annexitem}

\begin{annexitem}{annex:p08.klm.7}{thm:p08-klm-canonical-smooth}{An explicit minimal predecessor on every formula set}
\begin{proof}
Let \((m,\beta)\in\widehat\alpha\), so \(m\models\alpha\).
If \(\beta\leq\alpha\), the preceding characterization makes it
minimal. Otherwise the normal-world witness for
\(\alpha\vee\beta\not\sim\beta\) supplies
\(n\in N(\alpha\vee\beta)\) falsifying \(\beta\).
It satisfies \(\alpha\), and
\(\alpha\vee\beta\leq\beta\) makes
\((n,\alpha\vee\beta)\prec(m,\beta)\).
The additional identity \(\alpha\vee\beta\leq\alpha\), together
with \(n\models\alpha\), makes this predecessor minimal in
\(\widehat\alpha\). Thus the required minimal predecessor actually
exists for every nonminimal state in each definable set.
Neither a finite state space nor a globally well-founded order has been
used. Empty definable sets satisfy the condition vacuously.
\end{proof}
\end{annexitem}

\begin{annexitem}{annex:p08.klm.8}{thm:p08-klm-pref-representation}{Both consequence directions of the general representation}
\begin{proof}
Soundness proves that every smooth preferential model defines a P
relation. Conversely, start with an arbitrary P relation on the fixed
\((L,U)\). Its canonical state set is a subset of the set \(U\times L\),
its label map lands in \(U\), and its defined preference is a strict
partial order by item~\ref{annex:p08.klm.5}. Item~\ref{annex:p08.klm.7}
establishes smoothness on every definable set, so it is a preferential model.

If \(\alpha\sim\beta\) and \((m,\gamma)\) is minimal in
\(\widehat\alpha\), then \(m\models\alpha\) and
\(\gamma\leq\alpha\). The first normal-world transport makes
\(m\) \(\alpha\)-normal, and therefore \(m\models\beta\).
Every such minimal state satisfies the consequent, proving
\(\alpha\sim_W\beta\).
If instead \(\alpha\not\sim\beta\), the compactness witness
provides \(m\in N(\alpha)\) falsifying \(\beta\).
The state \((m,\alpha)\) is minimal in \(\widehat\alpha\) by
reflexivity of \(\leq\) and the minimal characterization. It is
therefore an actual counterstate to \(\alpha\sim_W\beta\).
This proves equality of the two relations in both directions, not just
preservation of assertions already in a finite base. The construction
uses all normal-world pairs, with no choice of one representative per
formula and no quotient that erases necessary state copies.
\end{proof}
\end{annexitem}

\begin{annexitem}{annex:p08.klm.9}{thm:p08-klm-kb-completeness}{Arbitrary conditional bases, with finite derivations}
\begin{proof}
Let \(R_K\) be the set of assertions having finite P derivations from
\(K\). It contains \(K\) and all Reflexivity instances. Each rule
has finitely many assertion premises; concatenate their finite
derivations and append the rule application. Consequently \(R_K\)
is P-closed. Fixed classical-validity premises of the sheet are external
side conditions, not extra conditional assertions to be placed in \(K\).
Induction on derivations shows that every P-closed relation containing
\(K\) contains \(R_K\); thus it is the least one.

Every model of \(K\) validates the rules by soundness, so satisfies
every assertion in \(R_K\). If \(\alpha\sim\beta\notin R_K\),
apply the general representation to \(R_K\). The canonical model
represents it exactly and therefore satisfies \(K\), while the normal
world witnessing this nonmembership yields a minimal
\(\alpha\)-state falsifying \(\beta\). Hence the assertion is
not true in every model of \(K\). This proves both completeness
directions for arbitrary set-sized \(K\). Its finitary derivations do
not make the canonical model finite or reduce the normal theory to a
finite one; compactness has already been used at its precise witness step.
\end{proof}
\end{annexitem}

\begin{annexitem}{annex:p08.klm.10}{thm:p08-klm-penguins}{All five source consequences and the exception model}
\begin{proof}
First, Cautious Monotonicity with \(p\sim b,p\sim\neg f\) gives
\(p\wedge b\sim\neg f\).
Or applied to \(p\sim b,b\sim b\) gives \(b\leq p\);
Or applied to \(b\sim f,f\sim f\) gives \(f\leq b\).
The first transport applied to \(b\leq p,p\sim\neg f\) yields
\(b\sim p\to\neg f\). And with \(b\sim f\), followed by
Right Weakening, gives \(b\sim\neg p\).
Preorder transitivity gives \(f\leq p\); the same transport and
\(f\sim f\) give \(f\sim\neg p\). These are particular
derived consequences, not a general Contraposition rule.
Let \(D=b\vee p\). We have \(D\sim b\) and \(b\sim D\),
the second by Supraclassicality. Mutual antecedent equivalence transfers
\(b\sim f\) and \(b\sim\neg p\) to \(D\), proving the last
two asserted consequences.

For the concrete model, take \(U\) to be all classical valuations on
\(b,p,f\) and label the two states as displayed. Its strict preference
has only \(s_0\prec s_1\). Every finite strict partial order is
smooth on every subset: descending through nonminimal predecessors must
terminate by finiteness and irreflexivity, and transitivity makes the
terminal minimal state a predecessor of the starting one. Hence the
model is smooth on all formula sets.
The \(p\)-set is \(\{s_1\}\), so \(p\sim_W b\) and
\(p\sim_W\neg f\) hold and \(p\sim_W f\) fails.
The \(b\)-set has both states, with minimum \(s_0\), so
\(b\sim_W f\) holds. The \(b\wedge p\)-set has only \(s_1\)
and falsifies \(f\), demonstrating nonmonotonicity under an additional
antecedent. The \(b\wedge f\)-set has only \(s_0\), with the
same minima as the \(b\)-set, so every consequent has the same truth
under both antecedents. The universe itself still permits penguin
valuations not satisfying \(b\); we did not replace the default
\(p\sim b\) by a hard constraint on \(U\).
\end{proof}
\end{annexitem}

\begin{annexitem}{annex:p08.klm.11}{thm:p08-klm-diamond}{Derivations and a model for each nonentailment}
\begin{proof}
For each assertion \(x\sim y\), S with antecedent \(\top\),
followed by Left Logical Equivalence, gives \(\top\sim x\to y\).
Apply this to the four members of \(K\), and combine them by And.
Their material implications entail \(\neg(p\wedge s)\), since
\(p\) gives \(e\) and \(s\) gives \(\neg e\); they also entail
\(\neg t\), since \(t\) gives both \(p\) and \(s\).
Right Weakening proves both source entailments.

For the employment failures use states \(u,v,w\), with
\(u\prec w,v\prec w\) and no other comparisons, and labels
\[
 \begin{array}{c|cccc}
 &t&p&s&e\\\hline
 u&0&1&0&1\\v&0&0&1&0\\w&1&1&1&\varepsilon
 \end{array}
 \quad(\varepsilon=0\text{ or }1).
\]
The minimal \(p\)-state is \(u\), minimal \(s\)-state is \(v\),
and the only \(t\)-state is \(w\). Thus both models satisfy every
assertion of \(K\). With \(\varepsilon=0\) the model refutes
\(t\sim e\); with \(\varepsilon=1\) it refutes
\(t\sim\neg e\). Both are finite smooth models, so soundness
establishes both claimed nonderivabilities.

For \(p\sim\neg s\), use a three-state chain \(v\prec u\prec w\),
including its transitive comparison, with labels
\(v=(0,0,1,0)\), \(u=(0,1,1,1)\), \(w=(1,1,1,1)\)
in the same coordinate order. The \(p\)-minimum \(u\) satisfies
\(e\) but also \(s\), the \(s\)-minimum \(v\) satisfies
\(\neg e\), and \(w\) satisfies the two \(t\)-defaults.
Thus \(K\) holds but the requested assertion fails.
For \(s\sim\neg p\), use \(u\prec v\prec w\) with labels
\(u=(0,1,0,1)\), \(v=(0,1,1,0)\), \(w=(1,1,1,0)\).
Its \(p\)-minimum is \(u\) and \(s\)-minimum is \(v\);
all defaults hold while the latter minimum falsifies \(\neg p\).

Finally extend the first model with \(\varepsilon=0\) by assigning
a fresh atom \(a\) false at \(u,v\) and true at \(w\).
The \(a\wedge p\)-set has just \(w\), which falsifies \(e\).
This model still satisfies \(K\), proving that irrelevant vocabulary
does not automatically justify that strengthening in bare P.
For each model, use the full unconstrained valuation universe on its
atoms; other valuations need not occur as state labels. All comparisons
and minimal sets have been specified, and finiteness gives smoothness
by the preceding finite-order argument.
\end{proof}
\end{annexitem}

\begin{annexitem}{annex:p08.klm.12}{thm:p08-klm-copies}{Why identifying the two labels changes the relation}
\begin{proof}
Let the three represented valuations be \(A=(1,0)\), \(B=(0,0)\),
\(C=(1,1)\); the missing valuation is \(D=(0,1)\).
The original model is finite and smooth. Its \(\top\)-minima have
labels \(A,B\), so \(\top\sim_W\neg q\).
For antecedent \(p\), both the \(A\)-state and the \(C\)-copy
above the \(B\)-state are minimal; therefore \(p\not\sim_W\neg q\).
For \(p\leftrightarrow q\), the \(B\)-state and the other
\(C\)-copy are minimal, so
\(p\leftrightarrow q\not\sim_W\neg q\).

Each valuation is definable by its two-literal characteristic formula
\(\chi_A,\chi_B,\chi_C,\chi_D\). In any smooth model,
\(\chi_X\sim\bot\) holds iff there is no state labeled \(X\):
if such a state exists, smoothness supplies a minimal one, and its
label cannot satisfy \(\bot\). The original model therefore forces
an equivalent model to have labels \(A,B,C\) and no \(D\).
If its label map is injective, it has exactly three states \(a,b,c\).
The failure of \(p\sim\neg q\) forces \(c\) to be minimal
in \(\{a,c\}\), so \(a\not\prec c\).
The failure of \(p\leftrightarrow q\sim\neg q\) forces it to
be minimal in \(\{b,c\}\), so \(b\not\prec c\).
Irreflexivity excludes \(c\prec c\), so \(c\) is minimal in
the entire state set. Since its label satisfies \(q\), this contradicts
\(\top\sim\neg q\). Hence no equivalent injectively labeled
smooth model exists. The proof uses assertions that identify both
label presence and local minimality, not an unsupported claim that
copies merely look necessary.
\end{proof}
\end{annexitem}

\begin{annexitem}{annex:p08.klm.13}{thm:p08-klm-nonsmooth}{Every definable set has minima but Cautious Monotonicity fails}
\begin{proof}
The displayed relation is a strict partial order: the descending-chain
component is transitive and irreflexive, the disjoint component has only
\(a\prec d\), and no mixed comparisons are present.
Formula sets over \(p,q\) are unions of the three represented label
classes \(A=\{a\}\), \(B=\{b_0\}\),
\(C=\{d,b_i:i\geq1\}\). If a nonempty formula set contains \(a\),
that state is minimal. If it omits \(a\) but contains class \(C\),
\(d\) is minimal, since its only predecessor is the omitted \(a\).
If neither condition holds, its only possible represented class is
\(B\), and \(b_0\) is minimal. Thus every nonempty definable set
has a minimal state, even though infinitely many states share class \(C\).

The only global minimal state is \(a\): \(d\) has predecessor
\(a\), and every \(b_i\) has predecessor \(b_{i+1}\).
No \(b_i\) lies above the minimal \(a\), so the \(\top\)-set
is not smooth. Its minima satisfy both \(p\) and \(q\), yielding
\(\top\sim_W p\) and \(\top\sim_W q\).
But \(\widehat p=\{a,b_0\}\), with both states minimal, and
\(b_0\) falsifies \(q\). Thus \(p\not\sim_W q\), violating
Cautious Monotonicity. The underlying universe can be the four ordinary
valuations, so compactness and classical semantics are not the failed
hypotheses; precisely smoothness is missing.
\end{proof}
\end{annexitem}

\begin{annexitem}{annex:p08.klm.14}{thm:p08-klm-not-rational}{The finite rational-monotony and ranking countercase}
\begin{proof}
Call the states \(u,v,w\), with \(u\prec w\) only. This is a
finite strict partial order, hence smooth. Its global minima are
\(u,v\), both satisfying \(q\); \(v\) satisfies \(p\), so not
all of them satisfy \(\neg p\). The \(p\)-set is \(\{v,w\}\),
with both states minimal, and \(w\) falsifies \(q\).
This gives the three asserted properties and refutes Rational Monotony
at antecedent \(\top\), added premise \(p\), consequent \(q\).
The induced relation nevertheless satisfies all of P by soundness.

If a total ranking represented this preference, incomparability of
\(u,v\) would require equal ranks, as would incomparability of
\(v,w\). But \(u\prec w\) requires strictly ordered ranks, a
contradiction. Thus smoothness here neither is nor entails ranking.
No converse ranked representation or rational-closure construction is
being inferred from this finite separating example.
\end{proof}
\end{annexitem}

\begin{annexitem}{annex:p08.klm.15}{thm:p08-klm-cumulative-separation}{The cumulative one-state countermodel to Or}
\begin{proof}
Let \(U\) be the two valuations on \(p\), and label a single state
\(s\) by the whole nonempty set \(U\), using the source's cumulative
satisfaction condition \(s\equiv\alpha\) iff every label-world
satisfies \(\alpha\). There are no preference comparisons. Every
definable state set is empty or \(\{s\}\), so smoothness holds.
Reflexivity and Left Logical Equivalence follow directly from the
satisfaction condition; Right Weakening follows by validity on every
world of the label. For Cut, if the antecedent is true at \(s\),
the first premise makes the intermediate consequent true there; its
conjunction with the antecedent is therefore true, and the second
premise gives the final consequent. If the antecedent is false, the
conclusion is vacuous. For Cautious Monotonicity, any state satisfying
the conjunction satisfies the antecedent, so its premise consequence
gives the consequent. This verifies all five native C rules directly.

Neither \(p\) nor \(\neg p\) is true at \(s\), because its two
label-worlds disagree. Each has an empty state set, so both assertions
to \(\bot\) hold vacuously. The tautology \(p\vee\neg p\) is
true at \(s\), whose label cannot satisfy \(\bot\), so its
assertion to \(\bot\) fails. Thus Or is false.
The reason is the failure of
\(\widehat{\alpha\vee\beta}=\widehat\alpha\cup\widehat\beta\)
for set-valued labels, not a failure of classical truth at each world.
This distinguishes the two source model types without substituting
a preferential model for the original cumulative presentation.
\end{proof}
\end{annexitem}
